\section{Введение}

Камеры времени пролёта (time-of-flight, ToF) позволяют получать измерения расстояния до наблюдаемой поверхности. Несколько таких камер могут использоваться для построения общего трёхмерного представления сцены. Однако карты глубины и облака точек первоначально заданы в отдельных системах координат. Для их объединения необходимо определить вращения и переносы между камерами, то есть выполнить внешнюю калибровку.

Эта задача отличается от внутренней калибровки, в которой оценивают параметры проекции и искажения изображения. Даже при известных внутренних параметрах неверная взаимная ориентация камер приводит к несовпадению наблюдений одной поверхности. При этом ошибки ToF-измерений могут зависеть от расстояния, свойств поверхности и условий получения данных~\cite{foix2011}. Поэтому оценку внешних параметров следует отделять от оценки качества самой глубины.

В данной работе рассматривается калибровка по объекту известной геометрии, например по кубу с различимыми элементами на гранях. Такой объект задаёт общую систему координат для наблюдений нескольких камер. Камеры могут видеть разные элементы мишени, если их координаты определены в одной системе и взаимное положение граней известно. Это позволяет исследовать конфигурации, для которых поиск соответствий непосредственно между изображениями сцены затруднён.

Цель исследования --- разработать и оценить способ определения взаимного положения ToF-камер по наблюдениям калибровочного объекта. В качестве базового решения рассматривается оценка позы по изображениям с последующим совместным уточнением. Дополнительное направление --- использование глубины и анализ чувствительности результата к погрешностям геометрии мишени. На текущем этапе сформулированы модель, требования к входным данным и план проверки; реализация и экспериментальная оценка остаются последующими этапами.
